\documentclass[10pt,a4paper]{amsart}
\usepackage[margin=19mm]{geometry}
\usepackage{amsmath,amssymb,mathtools}
\usepackage[scr=rsfs,cal=euler]{mathalfa}
\usepackage{xfrac}
\usepackage[unicode,hidelinks,bookmarks=false]{hyperref}
\newcommand{\ie}{i.e.\ }
\newcommand{\cf}{cf.\ }
\newcommand{\resp}{respectively\ }
\newcommand{\Subsection}{\subsection}
\providecommand{\nobreakdash}{\nobreak\hbox{-}}
\setlength{\emergencystretch}{2em}
\multlinegap0pt
\usepackage{bm}
\usepackage{bbm}
\usepackage{dsfont}
\usepackage{xspace}
\usepackage{cancel}
\usepackage{mathtools}
\usepackage{amsthm}
\newtheorem{theorem}{Theorem}
\newcommand{\U}{\operatorname{U}}
\newcommand{\Z}{\mathbb Z}
\title[Geometric Duality Between Constraints and Gauge Fields: Mirr]{Geometric Duality Between Constraints and Gauge Fields: Mirror Realization and Reduction Geometry on Principal Bundles (Theorem 6)}
\author{Dongzhe Zheng}
\date{Source: arXiv:2506.00728v3 (CC BY 4.0)}
\begin{document}
\maketitle

\noindent\emph{Source and licence.} Geometric Duality Between Constraints and Gauge Fields: Mirror Realization and Reduction Geometry on Principal Bundles. Version arXiv:2506.00728v3 (CC BY 4.0), distributed under the Creative Commons Attribution 4.0 International licence, as recorded in the arXiv licence metadata for this exact version and shown on the abstract page https://arxiv.org/abs/2506.00728v3. Full rights notice: licenses/arxiv-2506.00728-CC-BY-4.0.txt.\par\smallskip

\noindent\textit{Editorial scope.} Counted unit: the Theorem printed as Theorem 6 of the article (source0.tex lines 407--430, source label \texttt{thm:diagonalization-base}), delivered together with the article's own complete original proof of it (source0.tex lines 432--477, one \texttt{proof} environment of 46 lines). No mathematics is written, repaired or re-derived: statement and proof are the article's own text line for line. Required context retained uncounted: none. Every definition, display and label of those lines is retained unchanged. Omitted from this package and not redistributed: 1--406, 431--431, 478--1719. Nothing inside a retained excerpt is omitted or altered: every retained body is delivered exactly as the article's lines are written, so no omission receipt of any kind accompanies this package. Citations: the retained excerpts carry no citation command at all, so no bibliography entry of the article is transcribed and no reference is redistributed. Reference adaptation: none was needed and none was made; every reference inside the delivered unit resolves inside it, so no unresolved reference or citation is produced. Printed slips kept literal: no printed word, symbol, hypothesis or number of the counted statement or of its proof was altered. Layout and class: the article's own class is \texttt{amsart} with packages including fontenc, lmodern, amsmath, amssymb, amsthm, mathtools, mathrsfs, microtype, hyperref, geometry; the wrapper is the lane's pinned \texttt{amsart} editorial wrapper at 10pt on A4, which restates the theorem-like environments the retained text needs and adds the article's own macro definitions that the retained text needs (\texttt{\textbackslash U}, \texttt{\textbackslash Z}), with extra wrapper packages (bm, bbm, dsfont, xspace, cancel, mathtools). The delivered numbering is the unit's own. Assets: no retained span contains a figure environment, a graphics-inclusion command, a PDF-page inclusion, a table or a TikZ picture; no image file is copied into this package, the package declares no asset dependency and no image file is redistributed with it. Licence: version arXiv:2506.00728v3 of this article is distributed under the Creative Commons Attribution 4.0 International licence, as recorded in the arXiv licence metadata for this exact version and shown on the abstract page https://arxiv.org/abs/2506.00728v3. A full rights notice is delivered at licenses/arxiv-2506.00728-CC-BY-4.0.txt. Characters: every retained excerpt is pure ASCII, so the delivered engine, XeLaTeX with its default Latin Modern text font, reports no missing character.
\begin{theorem}[Topology and structure group of the diagonalization bundle]
\label{thm:diagonalization-base}
For $n\geq2$ there are natural homeomorphisms
\[
 \mathcal B_n\simeq(\mathcal F_n\times\mathcal F_n)/W
 \simeq(\U(n)\times\U(n))/K.
\]
Its fundamental group is $S_n$ and
$H^2(\mathcal B_n;\Z)\simeq\Z/2$.
The generator is the first Chern class of the complex line associated
to the sign representation of the covering
$\mathcal F_n\times\mathcal F_n\to\mathcal B_n$.

The normal subgroup $\Delta T$ is the kernel of the affine action
\[
 K\longrightarrow\operatorname{Aff}_W(T)=T\rtimes W,
 \qquad (n_1,n_2)\cdot z=n_2zn_1^{-1},
\]
and the displayed map induces $K/\Delta T\simeq T\rtimes W$.
Consequently $(\U(n)\times\U(n))/\Delta T\to\mathcal B_n$
is a principal $T\rtimes W$ bundle, and $\mathcal E_n$ is its
associated $T$ bundle under the affine action. For $n=1$ the base is a
point and its positive-degree cohomology vanishes.
\end{theorem}

\begin{proof}
Ordering the first projection family and transporting that order
through its bijection orders the second family. The $n!$ choices differ
by a simultaneous permutation, giving the first quotient. The source
is compact and $\mathcal B_n$ is Hausdorff in its metric topology, so
the induced continuous bijection is a homeomorphism. The second
quotient follows by lifting both ordered families to unitary frames.

The long exact sequence for $T\to\U(n)\to\mathcal F_n$ shows
$\pi_1(\mathcal F_n)=0$: the map
$\pi_1(T)=\Z^n\to\pi_1(\U(n))=\Z$ is summation and is surjective.
Thus $\mathcal F_n^2$ is the universal cover of $\mathcal B_n$,
and $H_1(\mathcal B_n;\Z)=(S_n)_{\mathrm{ab}}=\Z/2$.
The degree-two cohomology of the flag manifold is
\[
 H^2(\mathcal F_n;\mathbb Q)
 \simeq\mathbb Q^n/\mathbb Q(1,\ldots,1),
\]
with $S_n$ permuting coordinates. Its invariant subspace is zero for
$n\geq2$. Since $H^1(\mathcal F_n;\mathbb Q)=0$, the same holds for
$H^2(\mathcal F_n^2;\mathbb Q)^W$. Transfer for the finite covering
gives $H^2(\mathcal B_n;\mathbb Q)=0$. The base is a compact manifold
and has the homotopy type of a finite CW complex. Hence
$\operatorname{Hom}(H_2(\mathcal B_n;\Z),\Z)=0$, and the integral
universal coefficient sequence yields
\[
 H^2(\mathcal B_n;\Z)
 \simeq\operatorname{Ext}(H_1(\mathcal B_n;\Z),\Z)
 \simeq\Z/2.
\]
The Bockstein of the sign class in $H^1(\mathcal B_n;\Z/2)$ is nonzero
because $H^1(\mathcal B_n;\Z)=0$; it is the first Chern class stated
above.

The right $K$ action on $\U(n)^2$ is
$(u,v)\cdot(n_1,n_2)=(un_1,vn_2)$. If $z\in T$, the unitary operator
$vzu^{-1}$ maps the $i$th line of the first frame to its matched line
of the second. The associated-bundle relation
$[(un_1,vn_2),z]=[(u,v),n_2zn_1^{-1}]$ gives the model-fiber action
$z\mapsto n_2zn_1^{-1}$. Since $n_1,n_2$ have the same
Weyl component, $n_2n_1^{-1}\in T$ and this is an affine map of $T$.
It fixes every $z$ exactly when $n_1=n_2\in T$, proving the kernel
claim. Permutation matrices split $N\to W$, and every translation of
$T$ occurs, so the image is all $T\rtimes W$. The associated-bundle
description follows.
\end{proof}

\end{document}
